\documentclass[./main.tex]{subfiles}

\begin{document}
\chapter*{why this effort?}

Well, I contemplate about a lot of things and it seems that after a point I am not able to keep track of them.
I forget some important conclusions, perspectives, schools of thought, jewels of wisdom and such things in various aspects that are fruits of a lot of effort, suffering and pain.
And I often have to re-start my train of thought on them, which is tedious and dis-motivating after a point of time.

\paragraph{Therefore this journal is to}

\begin{enumerate}[noitemsep]
  \item Write down stuff
    \begin{enumerate}[noitemsep]
      \item In a simplified, clean, elegant, minimalistic form.
      \item So that I don't lose it permanently.
      \item So that I don't have to remember what I don't need to.
      \item So that I use my brain for thinking rather than storing.
    \end{enumerate}
  \item Use the power of machine search.
  \item Don't start over from scratch. Build on previous knowledge.
\end{enumerate}

\begin{figure}[h]
  \centering
  \begin{tikzpicture}
    \draw (0, 0) circle (1.5) node {Write Stuff};
    \draw (5, 0) circle (1.5) node {Search Stuff};
    \draw (10, 0) circle (1.5) node {Update Stuff};
  \end{tikzpicture}
\end{figure}

\paragraph{The effort will be to}
\begin{enumerate}[noitemsep]
  \item Build constructs that are simple, atomic, composable and excel at one thing.
  \item Build meaningful connections and flow among constructs.
  \item Keep everything as simple and minimalistic as possible.
  \item Promote recognition rather than re-call.
\end{enumerate}
These ideas are inspired by the UNIX philosophy, my experience in programming and doing things in general.

\end{document}
